\input{../../../templates/es/preambulo.tex}

% Los listados se toman de src/ con \lstinputlisting; la diapositiva es el archivo.
\lstset{
    basicstyle=\ttfamily\fontsize{8pt}{9.6pt}\selectfont,
    breaklines=true,
    breakatwhitespace=true,
    columns=fullflexible,
    keepspaces=true,
    xleftmargin=0.3em,
    framesep=2pt,
    aboveskip=0pt,
    belowskip=0pt
}

\newcommand{\marcoresultado}[3]{%
    \begin{frame}{#1}
        \begin{center}
            \includegraphics[height=0.62\textheight,width=\textwidth,keepaspectratio]{#2}
        \end{center}
        \vspace{0.25em}
        {\small #3\par}
    \end{frame}%
}

\newcommand{\marcoafirmacion}[2]{%
    \begin{frame}{#1}
        {\small #2\par}
    \end{frame}%
}

\date{}

\begin{document}

\tituloleccion{04}{Cruza}

\section{Esta lección}

\begin{frame}{Esta lección}
La Lección 03 abrió la selección sobre una generación. Sigue la cruza, todavía sobre dos padres fijos, con un solo instrumento.

\vspace{0.6em}
\begin{itemize}
    \item Una cruza se define por los hijos que puede alcanzar.
    \item Más puntos de corte no es un dial de fuerza.
    \item Cortar-y-copiar copia valores parentales; el blend calcula otros --- y se sale del intervalo.
\end{itemize}
\end{frame}

% ================= descendencia_01_dos_padres =================
\begin{frame}{El alcance es el soporte del operador}
Para BLX-\(\alpha\), sean \(m=\min(a,b)\), \(M=\max(a,b)\) y \(d=M-m\):
\[
x'\sim U[m-\alpha d,\,M+\alpha d],\qquad
\Pr(x'\notin[m,M])=\frac{2\alpha}{1+2\alpha}.
\]
La cruza uniforme tiene \(2^L\) primeros hijos posibles si todos los pares
parentales difieren; la lineal alcanza dos hijos afines. Los cortes canónicos de
(n) puntos usan (1,\ldots,L-1), pero este código excluye (L-1).
Muestrear ilustra el alcance; enumerar fija su cardinalidad. OX1 conserva
pertenencia y unicidad de una permutación.
\end{frame}

\section{Dos padres, y cómo medir lo que un operador les hace}

\begin{frame}{Dos padres, y cómo medir lo que un operador les hace}
La cruza es el operador que toma dos cromosomas y devuelve dos nuevos.
Nada más en un algoritmo genético crea una combinación que no estaba ahí
antes, por lo que toda la lección gira en torno a una pregunta: DADOS ESTOS DOS PADRES,
¿CUÁL ES EL CONJUNTO DE HIJOS QUE UN OPERADOR PUEDE PRODUCIR?
\end{frame}

\marcoresultado{Dos padres, y cómo medir lo que un operador les hace}{../figures/descendencia_01_dos_padres.png}{%
    La fila cero: 2 hijos alcanzados, 0.00\% genes nuevos, 100.00\% clones, 100.00\% legales. También el número contra el que se mide cada paso posterior: elegir libremente entre los padres permite 2⁶ = \textbf{64} cromosomas}

% ================= descendencia_02_un_punto =================
\section{Cruza de un punto, y el tamaño de su alcance}

\begin{frame}{Cruza de un punto, y el tamaño de su alcance}
El operador de libro de texto: cortar a ambos padres en el mismo punto, intercambiar las colas. Vale
la pena comenzar aquí no porque sea el mejor sino porque su conjunto alcanzable es lo
suficientemente pequeño como para escribirlo completo - que es la única manera de ver lo que el
operador puede y no puede hacer.
\end{frame}

\begin{frame}[fragile]{(1) cruza\_un\_punto()}
\lstinputlisting[basicstyle=\ttfamily\fontsize{8pt}{9.6pt}\selectfont,firstline=129,lastline=141]{../src/descendencia_02_un_punto.py}
\end{frame}

\begin{frame}[fragile]{(2) el conjunto alcanzable}
\lstinputlisting[basicstyle=\ttfamily\fontsize{6pt}{7.2pt}\selectfont,firstline=149,lastline=181]{../src/descendencia_02_un_punto.py}
\end{frame}

\begin{frame}[fragile]{(3) la figura de cortes}
\lstinputlisting[basicstyle=\ttfamily\fontsize{7pt}{8.5pt}\selectfont,firstline=184,lastline=209]{../src/descendencia_02_un_punto.py}
\end{frame}

\marcoresultado{Cruza de un punto, y el tamaño de su alcance}{../figures/descendencia_02_un_punto.png}{%
    2,000 extracciones producen \textbf{10} hijos distintos; enumerar los 5 puntos de corte produce los mismos 10. \textbf{0.00\%} genes nuevos: el operador elige qué padre, nunca qué valor}

\marcoresultado{El mecanismo: un corte, dos hijos}{../figures/descendencia_02_un_punto_esquema.png}{%
    Corte tras el gen 3. Los genes azules vienen del padre 1, los rojos del padre 2.
    El hijo es un padre hasta el corte y el otro después. Ningún gen se inventa.}

% ================= descendencia_03_mas_cortes =================
\section{Más puntos de corte, y lo que la familia nunca puede hacer}

\begin{frame}{Más puntos de corte, y lo que la familia nunca puede hacer}
Si un corte es bueno, ¿son tres mejor? Medido en este par la respuesta es no,
y la razón vale más que la respuesta. Lo que separa a los miembros de esta
familia no es cuántos cortes hacen sino cuántas de las 64 combinaciones de los
genes de los padres pueden realmente alcanzar: la cruza uniforme llega a todas,
los operadores de n puntos llegan a unas pocas.
\end{frame}

\begin{frame}[fragile]{(1) cruza\_n\_puntos()}
\lstinputlisting[basicstyle=\ttfamily\fontsize{7pt}{8.5pt}\selectfont,firstline=147,lastline=167]{../src/descendencia_03_mas_cortes.py}
\end{frame}

\begin{frame}[fragile]{(2) cruza\_uniforme()}
\lstinputlisting[basicstyle=\ttfamily\fontsize{8pt}{9.6pt}\selectfont,firstline=170,lastline=187]{../src/descendencia_03_mas_cortes.py}
\end{frame}

\begin{frame}[fragile]{(3) la tabla de alcance}
\lstinputlisting[basicstyle=\ttfamily\fontsize{6pt}{7.2pt}\selectfont,firstline=195,lastline=242]{../src/descendencia_03_mas_cortes.py}
\end{frame}

\begin{frame}[fragile]{(4) la figura de patrones}
\lstinputlisting[basicstyle=\ttfamily\fontsize{7pt}{8.5pt}\selectfont,firstline=245,lastline=267]{../src/descendencia_03_mas_cortes.py}
\end{frame}

\marcoresultado{Más puntos de corte, y lo que la familia nunca puede hacer}{../figures/descendencia_03_mas_cortes.png}{%
    El alcance no es monótono en los cortes: un punto \textbf{10}, 2 puntos \textbf{12}, 3 puntos \textbf{8}, uniforme \textbf{64 de 64}. El precio de la uniforme: \textbf{3.20\%} de sus hijos son un padre devuelto}

% ================= descendencia_04_mezcla =================
\section{Aritmética sobre genes, y los primeros hijos ilegales}

\begin{frame}{Aritmética sobre genes, y los primeros hijos ilegales}
Cada operador hasta ahora elegía entre los genes de los padres. Estos dos
calculan nuevos valores - lo cual solo es posible porque un gen aquí es un
NÚMERO, y los números pueden promediarse, interpolarse y extrapolarse. Esa
única suposición sobre la representación es lo que el paso 6 quitará.
\end{frame}

\begin{frame}[fragile]{(1) acotar()}
\lstinputlisting[basicstyle=\ttfamily\fontsize{8pt}{9.6pt}\selectfont,firstline=179,lastline=188]{../src/descendencia_04_mezcla.py}
\end{frame}

\begin{frame}[fragile]{(2) cruza\_mezcla()}
\lstinputlisting[basicstyle=\ttfamily\fontsize{6pt}{7.2pt}\selectfont,firstline=191,lastline=222]{../src/descendencia_04_mezcla.py}
\end{frame}

\begin{frame}[fragile]{(3) cruza\_lineal()}
\lstinputlisting[basicstyle=\ttfamily\fontsize{7pt}{8.5pt}\selectfont,firstline=225,lastline=246]{../src/descendencia_04_mezcla.py}
\end{frame}

\begin{frame}[fragile]{(4) la tabla de alcance}
\lstinputlisting[basicstyle=\ttfamily\fontsize{6pt}{7.2pt}\selectfont,firstline=254,lastline=340]{../src/descendencia_04_mezcla.py}
\end{frame}

\begin{frame}[fragile]{(5) la figura de mezcla}
\lstinputlisting[basicstyle=\ttfamily\fontsize{6pt}{7.2pt}\selectfont,firstline=343,lastline=372]{../src/descendencia_04_mezcla.py}
\end{frame}

\marcoresultado{Aritmética sobre genes, y los primeros hijos ilegales}{../figures/descendencia_04_mezcla.png}{%
    Los primeros genes calculados (\textbf{99.57\%} nuevos) y los primeros hijos ilegales. alpha = 0: \textbf{0.00\%} de genes más allá de los padres, 100\% legal — El colapso interno de la Lección 01, medido. alpha = 0.5: \textbf{49.63\%} más allá, solo \textbf{50.08\%} legal; con \texttt{acotar()}, 100\% legal y aún 49.63\% más allá. Lineal: 100.00\% genes nuevos y exactamente \textbf{2} hijos alcanzables}

% ================= descendencia_05_guiada_por_aptitud =================
\section{Dejar que el operador vea lo que produjo}

\begin{frame}{Dejar que el operador vea lo que produjo}
Cada operador hasta ahora era ciego: producía dos hijos y los entregaba
sin preguntar si eran buenos. La variante guiada por aptitud de Gridin
agrega una línea - clasifica los dos padres y los dos hijos juntos, devuelve
los mejores dos - y con ella la garantía de que la cruza nunca puede empeorar
el par.
\end{frame}

\begin{frame}[fragile]{(1) objetivo()}
\lstinputlisting[basicstyle=\ttfamily\fontsize{8pt}{9.6pt}\selectfont,firstline=182,lastline=191]{../src/descendencia_05_guiada_por_aptitud.py}
\end{frame}

\begin{frame}[fragile]{(2) cruza\_guiada\_por\_aptitud()}
\lstinputlisting[basicstyle=\ttfamily\fontsize{8pt}{9.6pt}\selectfont,firstline=194,lastline=213]{../src/descendencia_05_guiada_por_aptitud.py}
\end{frame}

\begin{frame}[fragile]{(3) los sobrevivientes}
\lstinputlisting[basicstyle=\ttfamily\fontsize{6pt}{7.2pt}\selectfont,firstline=221,lastline=290]{../src/descendencia_05_guiada_por_aptitud.py}
\end{frame}

\begin{frame}[fragile]{(4) la figura de aptitud}
\lstinputlisting[basicstyle=\ttfamily\fontsize{6pt}{7.2pt}\selectfont,firstline=293,lastline=324]{../src/descendencia_05_guiada_por_aptitud.py}
\end{frame}

\marcoresultado{Dejar que el operador vea lo que produjo}{../figures/descendencia_05_guiada_por_aptitud.png}{%
    La mezcla ciega es peor que sus padres en \textbf{41.70\%} de las extracciones — la cruza no garantiza nada. La regla guiada por aptitud es peor en \textbf{0.00\%}, y paga por ello: \textbf{16.50\%} de las extracciones devuelven a los padres sin cambios, y solo el \textbf{56.10\%} de lo que devuelve es un hijo}

% ================= descendencia_06_la_division =================
\section{Los mismos operadores, en un cromosoma que es una permutación}

\begin{frame}{Los mismos operadores, en un cromosoma que es una permutación}
Nada de los operadores cambia en este paso. El cromosoma cambia: nueve genes
en lugar de seis, y los genes son las nueve paradas de una ronda de entregas,
cada una de las cuales debe aparecer exactamente una vez.
\end{frame}

\begin{frame}[fragile]{(1) las nueve paradas}
\lstinputlisting[basicstyle=\ttfamily\fontsize{7pt}{8.5pt}\selectfont,firstline=196,lastline=217]{../src/descendencia_06_la_division.py}
\end{frame}

\begin{frame}[fragile]{(2) longitud\_recorrido()}
\lstinputlisting[basicstyle=\ttfamily\fontsize{8pt}{9.6pt}\selectfont,firstline=220,lastline=225]{../src/descendencia_06_la_division.py}
\end{frame}

\begin{frame}[fragile]{(3) es\_legal\_recorrido()}
\lstinputlisting[basicstyle=\ttfamily\fontsize{8pt}{9.6pt}\selectfont,firstline=228,lastline=244]{../src/descendencia_06_la_division.py}
\end{frame}

\begin{frame}[fragile]{(4) los mismos operadores}
\lstinputlisting[basicstyle=\ttfamily\fontsize{6pt}{7.2pt}\selectfont,firstline=252,lastline=328]{../src/descendencia_06_la_division.py}
\end{frame}

\begin{frame}[fragile]{(5) la figura del mapa}
\lstinputlisting[basicstyle=\ttfamily\fontsize{6pt}{7.2pt}\selectfont,firstline=331,lastline=364]{../src/descendencia_06_la_division.py}
\end{frame}

\marcoresultado{Los mismos operadores, en un cromosoma que es una permutación}{../figures/descendencia_06_la_division.png}{%
    Los mismos operadores, una ruta en lugar de un punto. \textbf{8 de 8} puntos de corte dan una ruta que visita una parada dos veces y otra ninguna. Hijos legales: un punto \textbf{0.00\%}, 3 puntos \textbf{0.00\%}, mezcla \textbf{0.00\%}, lineal \textbf{0.00\%}, uniforme \textbf{0.95\%} — y el recuento de hijos legales de la uniforme que no eran simplemente un padre es \textbf{0}}

% ================= descendencia_07_ordenada =================
\section{Cruza ordenada, y la división leída en ambas direcciones}

\begin{frame}{Cruza ordenada, y la división leída en ambas direcciones}
La cruza ordenada (OX1) es la respuesta al paso 6: nunca intercambia valores de genes
posición a posición, así que una permutación entra y una permutación sale, por
construcción en lugar de por suerte. Este paso verifica que es legal, verifica lo que
realmente hereda, y luego hace lo que cierra la lección - entrega OX1 los padres de
valores reales del paso 1 y observa qué pasa.
\end{frame}

\begin{frame}[fragile]{(1) cruza\_ordenada()}
\lstinputlisting[basicstyle=\ttfamily\fontsize{6pt}{7.2pt}\selectfont,firstline=233,lastline=263]{../src/descendencia_07_ordenada.py}
\end{frame}

\begin{frame}[fragile]{(2) la prueba de adyacencia}
\lstinputlisting[basicstyle=\ttfamily\fontsize{6pt}{7.2pt}\selectfont,firstline=271,lastline=329]{../src/descendencia_07_ordenada.py}
\end{frame}

\begin{frame}[fragile]{(3) la otra dirección}
\lstinputlisting[basicstyle=\ttfamily\fontsize{7pt}{8.5pt}\selectfont,firstline=332,lastline=355]{../src/descendencia_07_ordenada.py}
\end{frame}

\begin{frame}[fragile]{(4) la comparación completa}
\lstinputlisting[basicstyle=\ttfamily\fontsize{6pt}{7.2pt}\selectfont,firstline=358,lastline=425]{../src/descendencia_07_ordenada.py}
\end{frame}

\begin{frame}[fragile]{(5) la figura de la división}
\lstinputlisting[basicstyle=\ttfamily\fontsize{7pt}{8.5pt}\selectfont,firstline=428,lastline=449]{../src/descendencia_07_ordenada.py}
\end{frame}

\marcoresultado{Cruza ordenada, y la división leída en ambas direcciones}{../figures/descendencia_07_ordenada.png}{%
    OX1 es legal el \textbf{100.00\%} en las rutas y mantiene \textbf{7.17 de 9} tramos de los padres, frente a \textbf{3.71} de una ruta aleatoria. Al recibir los padres con valores reales no falla: alcanza \textbf{2} hijos, el \textbf{100.00\%} de ellos un padre. Tabla final: operadores útiles en ambas representaciones — \textbf{0}}

\section{Siguiente}

\begin{frame}{Siguiente}
En una ruta de nueve paradas, un punto, n puntos, blend y lineal producen un hijo legal el 0\% de las veces. La cruza ordenada es legal el 100\% en rutas y en genes reales devuelve a los padres en silencio. Cero operadores sirven en las dos.

\vspace{0.8em}
\textbf{Lección 05 --- Mutación}

El tercer operador. OX1 en un tour fue un no-op silencioso aquí; esa cuenta se paga en la Lección 10, donde el orden es el cromosoma.
\end{frame}
\end{document}
